% LSPV-5 NOTE::BOOLGRID1::2026-10-02T07:14:00-07:00
% !TEX program = xelatex
\documentclass[11pt]{article}\usepackage[margin=.88in]{geometry}\usepackage{amsmath,amssymb}
\begin{document}
\section*{BOOLGRID1: exact Boolean-grid quotient identity for vMS LDE leaves}
\textbf{STATUS: PROVED/EXACT source lemma; transfer to exterior growth OPEN.} P versus NP remains open.

\paragraph{Source setting.}
In the vMS PCPP branch, a canonical tensor-grid low-degree extension restricts on an axis-parallel grid line to
\[
f\in\mathbb F_p[t],\qquad \deg f\le h-1,
\]
with f(a) in {0,1} on the h distinct grid points a in H. The active parameters include h=m^{100}, Delta=hr and p at least Delta^{100}. The declared SU05/Hadamard orbit realization later evaluates f away from H and binary-labels those values.

\paragraph{PROVED/EXACT: Boolean-grid degree floor.}
Let H be any h distinct field elements and suppose f(a) in {0,1} for every a in H. If f is nonconstant, then
\[
\boxed{\deg f\ge \lceil h/2\rceil.}
\]
Proof: f(f-1) vanishes on all h points. If 2 deg f<h, root counting forces f(f-1) to be the zero polynomial. Since F_p[t] is an integral domain, f=0 or f=1, contradiction.

\paragraph{PROVED/EXACT: quotient identity.}
Let
\[
P_H(t)=\prod_{a\in H}(t-a).
\]
Then
\[
\boxed{f(t)^2-f(t)=P_H(t)q(t)}
\]
for a unique q in F_p[t]. If f is nonconstant then q is nonzero, and
\[
\deg q=2\deg f-h\le h-2.
\]
This is an exact algebraic signature of every nonconstant Boolean-grid LDE leaf.

\paragraph{Hostile audit.}
The COVERBREAK1 word 0000110 over F_13 with h=7 interpolates to degree 6=h-1, so the degree-floor lemma does not remove the known hostile finite family. Thus high odd-field degree alone cannot imply dim U=omega(m).

\paragraph{Potential orbit use and obstruction.}
Along an SU orbit x_t=A^t v_0,
\[
f(x_t)^2-f(x_t)=P_H(x_t)q(x_t).
\]
This gives a synchronized multiplicative identity before binary encoding. However the current Hadamard branch uses a binary label kappa:F_p -> F_2^k, and kappa does not preserve odd-field multiplication. Therefore no identity for the binary exterior span U follows without a new transfer theorem. Treating the quotient identity as a Boolean quadratic relation would be invalid.

\paragraph{Candidate next theorem.}
Search for a source-matched statement that combines (i) q nonzero, (ii) h=m^{100}, (iii) p enormous relative to h, and (iv) the multiplicative orbit to force many independent same-shift binary label pairs after Hadamard conversion. A weaker sufficient target remains
\[
\boxed{\dim U=\omega(m)}.
\]
An alternative is to prove that any low-exterior exception forces q into an explicit exceptional algebraic class, then test whether perfectly accepted PCPP leaves can belong to that class.

\paragraph{Strict ledger.}
Boolean-grid degree floor and quotient divisibility: PROVED/EXACT. COVERBREAK1 degree-6 check: COMPUTATIONAL exact finite check. Quotient-to-Hadamard exterior transfer: OPEN. Large-field asymptotic exterior theorem: OPEN. Acquisition, reconstruction, arbitrary-algorithm transfer, nonrelativizing/nonalgebrizing hinge and final contradiction: OPEN.

\begin{thebibliography}{2}
\bibitem{vMS25} D. van Melkebeek and N. Mocelin Sdroievski, Instance-Wise Hardness and Refutation versus Derandomization for Arthur--Merlin Protocols, Computational Complexity 34:15, 2025.
\bibitem{SU05} R. Shaltiel and C. Umans, Simple Extractors for All Min-Entropies and a New Pseudorandom Generator, J. ACM 52(2):172--216, 2005.
\end{thebibliography}
\end{document}
% END_OF_VERIFIED_NOTE::LSPV-5::BOOLGRID1